% R15 component; only input paths and enumerated location/grammar prose are edited.
% Source: revisions/2026-10-04-r15/manuscript/general_theory.tex; commit 1cb3cc9efe135966ec228dfaf51c6841a6dfc97c.
% Generated by code/integrate.py; edit extraction rules there.
% Reviewed source: ECTA.tex, line 181, commit f5021cefa71492babfcbfa580e0c984f59a9de26
\section{Verification, Approximation, and Viscosity Selection}\label{sec:theory}
For a feasible policy $a$ and a smooth candidate $v$, define
\begin{equation}\label{eq:errors}
 r^a[v]=-v_t-\HH^a[v],\qquad
 g^a[v]=\sup_{b\in\Act}\HH^b[v]-\HH^a[v]\geq0.
\end{equation}
The first error concerns evaluation; the second concerns improvement. Their difference is the maximized HJB residual: $F[v]=r^a[v]-g^a[v]$. We keep the two quantities separate because the policy value and its regret depend differently on them.

\begin{theorem}[Finite-horizon economic error bound]\label{thm:verification}
Under Assumption~\ref{ass:model}, suppose a bounded smooth candidate $v$ in the declared utility interval and an admissible Markov policy $a$ satisfy $|r^a[v]|\leq e$, $0\leq g^a[v]\leq\eta$ throughout the verification domain. Suppose the boundary payoff error is at most $B$. Any reflected boundary condition is imposed exactly. Set $A_L(s)=(e^{Ls}-1)/L$ for $L>0$ and $A_0(s)=s$. Then, for $s=T-t$,
\begin{align}
 |v(t,x)-V^a(t,x)|&\leq e^{Ls}B+A_L(s)e,\label{eq:policybound}\\
 |v(t,x)-V^*(t,x)|&\leq e^{Ls}B+A_L(s)(e+\eta),\label{eq:valuebound}\\
 0\leq V^*(t,x)-V^a(t,x)&\leq2e^{Ls}B+A_L(s)(2e+\eta).\label{eq:regret}
\end{align}
Consequently, exact evaluation, exact boundary data, and a zero feasible-action gap identify the optimal value and an optimal policy.
\end{theorem}

The theorem replaces an assertion about a global minimizer of a composite loss with a statement about the exact computational objects NBO estimates. It allows errors, does not require that a finite network class contain $V^*$, and bounds performance against all admissible policies covered by comparison. A small sampled residual is not yet the uniform hypothesis of the theorem. Section~\ref{sec:coverage} addresses that distinction.

\subsection{A positive-weight recursive Bellman approximation}
Let $P_h^a$ be a nonnegative transition/interpolation operator on a finite state set, with row sums one after adjoining absorbing states. Boundary values are fixed. For a continuation array $w$, define $\Psi_h^a w$ pointwise as the scalar solution
\begin{equation}\label{eq:implicit}
 z=P_h^a w+h f(x,a,z),\qquad
 T_h w=\max_{a\in\Act_h(x)}\Psi_h^a w.
\end{equation}
For stationary calculations assume $f_y\leq-\lambda<0$ on an invariant value domain, and require that all scalar roots remain in that domain. The map $z\mapsto z-hf(x,a,z)$ is strictly increasing with slope at least $1+\lambda h$. Hence the root is unique and $\Psi_h^a$ is order preserving and a sup-norm contraction with factor $q_h=(1+\lambda h)^{-1}$. The maximum over feasible actions preserves these properties. A root bracket must contain the root; arbitrary clipping of utility values is not part of the definition.

For a finite horizon, the same implicit construction is used backward from the prescribed terminal array. Under the Lipschitz condition in Assumption~\ref{ass:model}, $hL<1$ gives a one-step Lipschitz constant at most $(1-hL)^{-1}$. Discount contraction is not assumed for a recursive parameter regime where it has not been verified. For a time-additive problem, the positive-weight backup $hr_0+e^{-\rho h}P_h^a w$ is also consistent and monotone. It is the backup used by the NDU reference laboratory.

\begin{theorem}[Discrete NBO certificate]\label{thm:discrete}
Suppose \eqref{eq:implicit} is well defined on an invariant domain, with contraction factor $q_h<1$. Let $v_h^*$ and $v_h^a$ denote the fixed points of $T_h$ and $T_h^a$, respectively. For any neural value array $v$ and deterministic feasible policy $a$, set
\[
 \delta=\|v-T_h^av\|_\infty,\qquad
 \epsilon=\|T_hv-T_h^av\|_\infty.
\]
Then
\begin{equation}\label{eq:discretebound}
 \|v-v_h^a\|_\infty\leq\frac\delta{1-q_h},\quad
 \|v-v_h^*\|_\infty\leq\frac{\delta+\epsilon}{1-q_h},\quad
 \|v_h^*-v_h^a\|_\infty\leq\frac{2\delta+\epsilon}{1-q_h}.
\end{equation}
The certificate covers the stated finite state and action model. Continuous actions require an additional certified maximization or action-coverage error.
\end{theorem}

This result applies regardless of how the neural array was obtained. Neural projection need not itself be monotone. What matters is that the output satisfies the monotone Bellman equation to the displayed tolerance. The distinction makes it possible to use flexible representations while retaining a comparison-based acceptance test.

\begin{theorem}[Viscosity-consistent neural Bellman approximation]\label{thm:viscosity}
Consider stationary problems satisfying comparison in a bounded continuous class and the strict utility monotonicity used in Theorem~\ref{thm:discrete}. Suppose the positive-weight schemes are stable, locally consistent with the HJB operator, and attain the prescribed boundary data in the limit. The feasible action approximations become dense with vanishing maximization error. If neural arrays and feasible actors satisfy
\begin{equation}\label{eq:scalederror}
 \delta_h+\epsilon_h=o(h),
\end{equation}
then their interpolated values converge locally uniformly to the unique viscosity solution. Uniform convergence up to the boundary follows when boundary barriers give uniform boundary attainment. Their discrete policy-value arrays converge to the same limit. No assertion of convergence of the policy functions themselves is required.
\end{theorem}

Here stability, local consistency, comparison, and boundary attainment are numerical-analysis hypotheses, not consequences of calling a representation neural. The proof is given in the appendix. A finite-horizon version uses backward error accumulation: terminal error tends to zero and the sum of one-step errors tends to zero after multiplication by the one-step Lipschitz factors. This covers finite-horizon recursive problems without imposing stationary discount contraction.

For a diffusion with $m$ Brownian factors, an interior semi-Lagrangian construction uses the $2m$ locations $x+h b\pm\sqrt{mh}\,\sigma e_j$, each with weight $1/(2m)$. Positive multilinear interpolation preserves monotonicity. For smooth test functions its interpolation contribution is $O(\Delta x^2/h)$; thus $\Delta x^2/h\to0$ is a sufficient interior coupling condition. The mean and covariance of the cubature increments match those of the specified Brownian diffusion to first order. Stopping and reflection require consistent boundary transitions. Negative interpolation weights are not permitted in this construction.

Appendix~\ref{app:r15evaluation} gives the differential-residual
counterexample, independent domain and action coverage, the complete-action
envelope, and the guarded finite-model evaluation map. Appendix~\ref{app:proofs}
proves the results, including finite-horizon error accumulation.
